\documentclass[11pt,a4paper]{article}
\usepackage[utf8]{inputenc}
\usepackage[T1]{fontenc}
\usepackage[portuguese]{babel}
\usepackage{geometry}
\geometry{margin=2.5cm}
\usepackage{listings}
\usepackage{xcolor}
\usepackage{amsmath}
\usepackage{amssymb}
\usepackage{hyperref}
\usepackage{enumitem}
\usepackage{tcolorbox}
\usepackage{tabularx}
\usepackage{booktabs}

\definecolor{codebg}{rgb}{0.95,0.95,0.95}
\definecolor{codegreen}{rgb}{0.1,0.5,0.1}
\definecolor{codepurple}{rgb}{0.5,0.1,0.5}
\definecolor{codegray}{rgb}{0.5,0.5,0.5}

\lstdefinestyle{python}{
    backgroundcolor=\color{codebg},
    commentstyle=\color{codegray}\itshape,
    keywordstyle=\color{codepurple}\bfseries,
    stringstyle=\color{codegreen},
    basicstyle=\ttfamily\small,
    breaklines=true,
    frame=single,
    language=Python,
    showstringspaces=false,
    tabsize=4,
    literate={á}{{\'a}}1 {ã}{{\~a}}1 {é}{{\'e}}1 {í}{{\'i}}1 {ó}{{\'o}}1 {õ}{{\~o}}1 {ú}{{\'u}}1 {ç}{{\c{c}}}1 {Á}{{\'A}}1 {Ã}{{\~A}}1 {É}{{\'E}}1 {Í}{{\'I}}1 {Ó}{{\'O}}1 {Õ}{{\~O}}1 {Ú}{{\'U}}1 {Ç}{{\c{C}}}1 {â}{{\^a}}1 {ê}{{\^e}}1 {ô}{{\^o}}1 {û}{{\^u}}1 {Â}{{\^A}}1 {Ê}{{\^E}}1 {Ô}{{\^O}}1 {Û}{{\^U}}1 {à}{{\`a}}1 {è}{{\`e}}1 {ì}{{\`i}}1 {ò}{{\`o}}1 {ù}{{\`u}}1 {ä}{{\"a}}1 {ö}{{\"o}}1 {Ä}{{\"A}}1 {Ö}{{\"O}}1
}
\lstset{style=python}

\newtcolorbox{objetivos}[1][]{
  colback=green!5!white,
  colframe=green!40!black,
  title=O que você vai aprender nesta página,
  fonttitle=\bfseries,
  #1
}

\newtcolorbox{exercicio}[1]{
  colback=blue!5!white,
  colframe=blue!40!black,
  title=#1,
  fonttitle=\bfseries
}

\newtcolorbox{solucao}{
  colback=yellow!5!white,
  colframe=orange!60!black,
  title=Solução proposta,
  fonttitle=\bfseries
}

\title{\textbf{Data Analysis with Python 2024--2025}\\Parte 6: Machine Learning Types\\\large Soluções dos Exercícios}
\author{Tradução e Adaptação: University of Helsinki --- MOOC.fi}
\date{}

\begin{document}

\maketitle

\section*{Nota Importante}
Este material é uma tradução e adaptação baseada no curso \textit{Data Analysis with Python 2024-2025}, oferecido pela \textbf{Universidade de Helsinque (MOOC.fi)}. As soluções apresentadas a seguir foram elaboradas para fins didáticos, seguindo os enunciados dos exercícios propostos no curso original.

\tableofcontents
\newpage

\section{Soluções - Capítulo 6}

\subsection*{Exercício 6.1 -- Classificação de Blobs}

\begin{solucao}
\begin{lstlisting}
from sklearn.naive_bayes import GaussianNB
from sklearn.model_selection import train_test_split
from sklearn.metrics import accuracy_score

def blob_classification(X, y):
    # Divide os dados de maneira deterministica (random_state=0) com tamanho de treino 75%
    X_train, X_test, y_train, y_test = train_test_split(X, y, train_size=0.75, random_state=0)
    
    # Treina o modelo usando Naive Bayes Gaussiano
    model = GaussianNB()
    model.fit(X_train, y_train)
    
    # Efetua as previsoes e retorna a acuracia
    predictions = model.predict(X_test)
    return accuracy_score(y_test, predictions)
\end{lstlisting}
\end{solucao}

\subsection*{Exercício 6.2 -- Classificação de Plantas}

\begin{solucao}
\begin{lstlisting}
from sklearn.datasets import load_iris
from sklearn.naive_bayes import GaussianNB
from sklearn.model_selection import train_test_split
from sklearn.metrics import accuracy_score

def plant_classification():
    # Carrega o conjunto local e divide com base em 80% do traning set
    dataset = load_iris()
    X = dataset.data
    y = dataset.target
    
    X_train, X_test, y_train, y_test = train_test_split(X, y, train_size=0.80, random_state=0)
    
    # Aplica o modelo Gaussiano
    model = GaussianNB()
    model.fit(X_train, y_train)
    
    # Retorna o score final
    return accuracy_score(y_test, model.predict(X_test))
\end{lstlisting}
\end{solucao}

\subsection*{Exercício 6.3 -- Classificação de Palavras}

\begin{solucao}
\begin{lstlisting}
import numpy as np
from sklearn.naive_bayes import MultinomialNB
from sklearn.model_selection import cross_val_score, KFold

ALPHABET = "abcdefghijklmnopqrstuvwxyzäö-"

def get_features(a):
    # a: array 1D com palavras
    features = np.zeros((len(a), len(ALPHABET)))
    for i, word in enumerate(a):
        for char in word:
            if char in ALPHABET:
                features[i, ALPHABET.index(char)] += 1
    return features

def contains_valid_chars(s):
    # Verifica validade exclusiva pelas letras do alfabeto permitido
    return all(c in ALPHABET for c in s)

def get_features_and_labels():
    # Simulando as cargas de listas locais:
    # finnish = load_finnish()
    # english = load_english()
    
    # (Para uso de solucao online, consideraremos lists importadas)
    
    # Processa finlandes
    finnish_filtered = []
    for w in finnish:
        w_lower = w.lower()
        if contains_valid_chars(w_lower):
            finnish_filtered.append(w_lower)
            
    # Processa ingles (Removendo nomes proprios / Letra maiuscula primeiro)
    english_filtered = []
    for w in english:
        if not w[0].isupper():
            w_lower = w.lower()
            if contains_valid_chars(w_lower):
                english_filtered.append(w_lower)
                
    X = get_features(np.array(finnish_filtered + english_filtered))
    y = np.array([0]*len(finnish_filtered) + [1]*len(english_filtered))
    return X, y

def word_classification():
    X, y = get_features_and_labels()
    model = MultinomialNB()
    
    # Usa cross-validation kfold para prever de forma fragmentada (em 5 iteracoes)
    cv = KFold(n_splits=5, shuffle=True, random_state=0)
    scores = cross_val_score(model, X, y, cv=cv)
    return scores
\end{lstlisting}
\end{solucao}

\subsection*{Exercício 6.4 -- Detecção de Spam}

\begin{solucao}
\begin{lstlisting}
import gzip
import numpy as np
from sklearn.feature_extraction.text import CountVectorizer
from sklearn.naive_bayes import MultinomialNB
from sklearn.model_selection import train_test_split
from sklearn.metrics import accuracy_score

def spam_detection(random_state=0, fraction=1.0):
    # Le os textos do ambiente src e captura as linhas referentes as fracoes processadas
    with gzip.open('src/ham.txt.gz', 'rt', encoding='utf-8') as f:
        ham_lines = f.readlines()
        ham = ham_lines[:int(fraction * len(ham_lines))]
        
    with gzip.open('src/spam.txt.gz', 'rt', encoding='utf-8') as f:
        spam_lines = f.readlines()
        spam = spam_lines[:int(fraction * len(spam_lines))]
        
    X = ham + spam
    y = np.array([0]*len(ham) + [1]*len(spam))
    
    # Inicializa a maquina de pipeline que converte letras repetidas por caracteristica
    vec = CountVectorizer()
    X_features = vec.fit_transform(X)
    
    # Divide dados
    X_train, X_test, y_train, y_test = train_test_split(X_features, y, train_size=0.75, random_state=random_state)
    
    # Classifica
    model = MultinomialNB()
    model.fit(X_train, y_train)
    predictions = model.predict(X_test)
    
    # Resultados
    acc = accuracy_score(y_test, predictions)
    total_tested = len(y_test)
    misclassified = (y_test != predictions).sum()
    
    return acc, total_tested, misclassified
\end{lstlisting}
\end{solucao}

\end{document}
